% ============================================================================
% RASCUNHO COMENTADO: Conclusão
% Para publicar, remova o "%" do início de cada linha.
% ============================================================================
\chapter{Conclusão}
\label{cap:conclusao}

% O trabalho partiu de um problema prático: quando a telemetria de um
% inversor fotovoltaico falha, a potência gerada deixa de ser registrada, e
% a lacuna compromete o acompanhamento da usina. A proposta foi reconstruir
% essa potência a partir de variáveis meteorológicas medidas de forma
% independente, com um soft-sensor embarcado, e avaliar o resultado tanto no
% computador quanto no hardware que seria usado em campo.
%
% O pipeline cobriu a preparação dos dados em quatro granularidades, o
% treino do Random Forest e do MLP, a geração das bibliotecas embarcadas, o
% firmware e o ensaio na placa. A comparação controlada mostrou que o MLP
% tem o menor erro em todas as granularidades e que o Random Forest fica
% próximo, com diferença pequena em favor da rede neural. O Random Forest
% leva vantagem no custo embarcado, por ser código \textit{inline} sem arena
% de memória. A regressão linear de referência fica bem abaixo, sobretudo em
% 1~d, com R$^2 = 0{,}534$.
%
% A análise cruzada de granularidades mostrou que a resolução de treino não
% prende o modelo à resolução de operação. Um modelo treinado em 1~h e
% avaliado em 5~min mantém R$^2$ de $0{,}784$ (MLP) e $0{,}777$ (RF); o
% modelo treinado em 1~d perde mais ao descer para o dado fino, com
% $0{,}696$ e $0{,}640$, mas se recupera em 1~h. Isso dá flexibilidade ao
% projeto, ao custo de uma perda que cresce com a distância entre as
% granularidades.
%
% O ensaio E3, com o Random Forest de 1~h na ESP32-S3, percorreu as 4\,042
% linhas do teste de 2025 do inversor E1. Fora dos \textit{gaps}, o caminho
% embarcado reproduziu o modelo treinado, com MAE de $0{,}066$~kW, RMSE de
% $0{,}083$~kW, R$^2$ de $0{,}972$ e viés de $+0{,}003$~kW. Dentro dos
% \textit{gaps}, o MAE foi de $0{,}065$~kW, enquanto a interpolação linear
% errou $0{,}312$~kW e a persistência, $0{,}376$~kW. Os números mostram a
% viabilidade do soft-sensor como complemento ao medidor: o erro dentro dos
% \textit{gaps} fica próximo do erro normal e muito abaixo das alternativas
% clássicas, que pioram à medida que o buraco cresce.
%
% As contribuições do trabalho são o pipeline completo, do dado bruto ao
% ensaio embarcado; a comparação controlada entre um modelo de árvores e uma
% rede neural nas mesmas condições; a análise cruzada de granularidades; o
% ensaio com \textit{gaps} artificiais e baselines de imputação; e a
% automação da geração dos artefatos no build.
%
% O trabalho tem limitações claras. Apenas o ensaio E3 foi executado; os
% ensaios E1, E2 e E4 a E8 e os inversores E2 e E3 permanecem pendentes. O
% ensaio é \textit{offline}: o firmware recebe lotes do servidor em vez de
% ler sensores. Os \textit{gaps} são artificiais, não falhas reais de
% telemetria. A calibração usa apenas o segundo semestre de 2024, e o
% estudo cobre um \textit{campus} e inversores de 2~kW. Essas limitações
% delimitam o alcance dos resultados, mas não mudam o que foi observado no
% ensaio executado.
%
% Os desdobramentos naturais são completar a matriz experimental e ensaiar
% os três inversores; levar o firmware para campo, com os sensores reais;
% comparar outras arquiteturas e outras formas de quantização; e integrar o
% soft-sensor ao sistema de aquisição da usina, para que a estimativa seja
% produzida no momento da falha e registrada junto com a série medida.
